% titel: Blatt 1, A4: Summe gerader Zahlen per Induktion

% Teil I: Werkzeuge, kein Lösungsweg.

Beweis per vollständiger Induktion, dass die Summe der ersten $n$ geraden
Zahlen $n(n+1)$ ist.

\begin{werkzeugkasten}[Vollständige Induktion]
Eine Aussage $A(n)$ für alle $n\ge1$:
\begin{enumerate}
  \item \textbf{Induktionsanfang} (IA): $A(1)$ prüfen.
  \item \textbf{Induktionsschritt} (IS): unter der Induktionsannahme
        $A(n)$ gilt $A(n+1)$.
\end{enumerate}
Dann gilt $A(n)$ für alle $n\ge1$.
\end{werkzeugkasten}

\begin{werkzeugkasten}[Summe gerader Zahlen]
\begin{equation}
  2+4+6+\dots+2n = 2(1+2+\dots+n)=2\cdot\frac{n(n+1)}2=n(n+1).
\end{equation}
Im Induktionsschritt zerlegt man
\begin{equation}
  \text{(Summe bis } n+1)=\text{(Summe bis } n)+ 2(n+1)
\end{equation}
und setzt die Induktionsannahme ein.
\end{werkzeugkasten}

\begin{achtung}[Darauf achten]
\begin{itemize}
  \item Im IS den \emph{neuen} Summanden $2(n+1)$ korrekt ans Ende der
        Summe stellen (nicht mit dem $n$-ten Summanden $2n$ verwechseln).
  \item Den gemeinsamen Faktor $2$ vor dem Klammern sum auf den neuen
        Summanden abschätzen: $n(n+1)+2(n+1)=(n+1)(n+2)$.
  \item Anfang bei $n=1$: $2=1\cdot2$.
\end{itemize}
\end{achtung}
